\begin{lemma}[Matrix product vector of the diagonal tensor]
    \label{lem:mpv_diagonal_tensor}
    \lean{MPOTensor.mpv_diagonalTensor}
    \leanok
    \uses{def:mpo_to_mps}
    For an MPO tensor $M$ and a configuration $\sigma$, the matrix product vector of
    the diagonal tensor $i\mapsto M^{ii}$ equals the matrix product vector of the
    doubled-index tensor $\widehat{M}$ evaluated on the diagonal-paired configuration
    $k\mapsto(\sigma_k,\sigma_k)$:
    $V^{(N)}(i\mapsto M^{ii})(\sigma)
    =V^{(N)}(\widehat{M})(k\mapsto(\sigma_k,\sigma_k))$.
\end{lemma}
\begin{proof}\leanok
    At each site value $i$, the diagonal tensor agrees with the doubled-index tensor
    on the diagonal pair, $M^{ii}=\widehat{M}^{\,(i,i)}$.
    The matrix product vector is multiplicative over sites, so replacing each factor
    $M^{\sigma_k\sigma_k}$ by $\widehat{M}^{\,(\sigma_k,\sigma_k)}$ reindexes the
    configuration from $\sigma$ to $k\mapsto(\sigma_k,\sigma_k)$, which is the stated
    equality.
\end{proof}

\begin{lemma}[Diagonal tensor as a block decomposition under horizontal form]
    \label{lem:mpv_diagonal_tensor_eq_blocks}
    \lean{MPOTensor.mpv_diagonalTensor_eq_blocks}
    \leanok
    \uses{lem:mpv_diagonal_tensor, def:horizontal_cf_mpdo}
    Suppose the doubled-index tensor $\widehat{M}$ shares its positive-length matrix
    product vector family with a block-diagonal direct sum
    $\bigoplus_k\mu_kA_k$ of blocks $A_k$ on
    the doubled index, as in a horizontal canonical form. Then the diagonal tensor of
    $M$ generates the same positive-length matrix product vector family as the
    block-diagonal direct sum
    $\bigoplus_k\mu_kA_k^{\mathrm{diag}}$ on the single-spin index, where
    $A_k^{\mathrm{diag}}(i)=A_k(i,i)$ is the diagonal restriction of each block.
    For every $N\geq1$,
    \begin{align}
        V^{(N)}(i\mapsto M^{ii})(\sigma)
        &=V^{(N)}\Bigl(\textstyle\bigoplus_k
          \mu_kA_k^{\mathrm{diag}}\Bigr)(\sigma).
        \notag
    \end{align}
\end{lemma}
\begin{proof}\leanok
    \uses{lem:mpv_diagonal_tensor}
    By Lemma~\ref{lem:mpv_diagonal_tensor}, the diagonal matrix product vector equals
    that of $\widehat{M}$ on the diagonal-paired configuration, which by hypothesis
    equals that of the block-diagonal direct sum there. Expanding the direct sum as the
    weighted sum $\sum_k\mu_k^NV^{(N)}(A_k)$ and restricting each block to the
    diagonal pair gives $\sum_k\mu_k^NV^{(N)}(A_k^{\mathrm{diag}})$, the matrix
    product vector of the single-spin block-diagonal direct sum.
\end{proof}

\begin{theorem}[Per-copy horizontal hypotheses do not control diagonal normality]
    \label{thm:diagonal_restriction_not_normal}
    \lean{MPSTensor.horizontalCFData_diagBlock_not_isNormal}
    \leanok
    \uses{rem:word_tuple_span_top}
    There is a left-canonical, injective doubled-index tensor with a nonzero
    block weight satisfying the older per-copy finite-length separation
    hypothesis whose restriction to the diagonal physical pairs is not normal.
    On bond space $\C^2$, take
    \begin{align}
        A^{(a,b)}&=c_aE_{ab},
        &
        c_0&=\frac35,
        &
        c_1&=\frac45.
        \notag
    \end{align}
    Then the four matrices $A^{(a,b)}$ span $\MN{2}$ and
    $\sum_{a,b}(A^{(a,b)})^*A^{(a,b)}=\Id$.
    For a single block, injectivity also gives finite-length block separation.
    The diagonal restriction, however, consists only of
    \begin{align}
        A^{\mathrm{diag}}(0)&=\frac35E_{00},
        &
        A^{\mathrm{diag}}(1)&=\frac45E_{11}.
        \notag
    \end{align}
    Every product of these two matrices is diagonal, so no blocking length spans
    $\MN{2}$.
\end{theorem}
\begin{proof}\leanok
    The four nonzero scalar multiples of the matrix units span the full matrix
    algebra. The identity $(3/5)^2+(4/5)^2=1$ gives left-canonicality. After
    diagonal restriction, $E_{ii}E_{jj}=\delta_{ij}E_{ii}$.
    Hence every word is either zero or a scalar multiple of some $E_{ii}$, and
    therefore has vanishing $(0,1)$ entry. Since $E_{01}$ does not, it lies
    outside the span at every blocking length.
\end{proof}

\begin{lemma}[Diagonal tensor word evaluation]
    \label{lem:eval_word_diagonal_tensor}
    \lean{MPOTensor.evalWord_diagonalTensor}
    \leanok
    \uses{def:mpo_eval_word}
    For any word $w$, the matrix product of the diagonal tensor along $w$ equals the
    operator word evaluation with equal ket and bra words, since both multiply the
    matrices $M^{w_kw_k}$ site by site.
\end{lemma}
\begin{proof}\leanok
    By induction on the word.
\end{proof}

\begin{lemma}[Diagonal tensor evaluates the density-operator diagonal]
    \label{lem:mpv_diagonal_tensor_eq_density}
    \lean{MPOTensor.mpv_diagonalTensor_eq_mpo_diag}
    \leanok
    \uses{def:mpo_family, lem:eval_word_diagonal_tensor}
    For an MPO tensor $M$ and a configuration $\sigma$, the matrix product vector of
    the diagonal tensor equals the diagonal entry of the generated operator:
    \begin{align}
        V^{(N)}(i\mapsto M^{ii})(\sigma)
        &=\bra{\sigma}\rho^{(N)}(M)\ket{\sigma}.
        \label{eq:mpdo_diagonal_density}
    \end{align}
    Both sides of (\ref{eq:mpdo_diagonal_density}) equal
    $\tr(M^{\sigma_1\sigma_1}\cdots M^{\sigma_N\sigma_N})$.
\end{lemma}
\begin{proof}\leanok
    \uses{lem:eval_word_diagonal_tensor}
    Word evaluation of the diagonal tensor multiplies the matrices
    $M^{\sigma_k\sigma_k}$, which is exactly the operator word evaluation with equal
    ket and bra indices; taking the trace gives the diagonal entry.
\end{proof}

\begin{lemma}[Nonnegativity of the diagonal matrix product vector]
    \label{lem:mpv_diagonal_tensor_nonneg}
    \lean{MPOTensor.mpv_diagonalTensor_nonneg}
    \leanok
    \uses{lem:mpv_diagonal_tensor_eq_density, def:mpdo}
    If $\rho^{(N)}(M)$ is positive semidefinite, then
    $V^{(N)}(i\mapsto M^{ii})(\sigma)\geq0$ for every configuration
    $\sigma$.
\end{lemma}
\begin{proof}\leanok
    By Lemma~\ref{lem:mpv_diagonal_tensor_eq_density}, the value is a diagonal
    entry of $\rho^{(N)}(M)$, and the diagonal entries of a positive semidefinite
    matrix are non-negative.
\end{proof}

\begin{lemma}[Matrix product vector of the vertical direct-sum tensor]
    \label{lem:mpv_vertical_assembled}
    \lean{MPOTensor.verticalCopyDim}
    \lean{MPOTensor.verticalCopyBlocks}
    \lean{MPOTensor.mpv_verticalAssembledTensor_eq_sum}
    \leanok
    \uses{thm:mpv_decomposition, def:block_diagonal_tensor}
    Let $g$ be a block count,
    $\dim,\mult:\{0,\ldots,g-1\}\to\N$, weights
    $\omega_{\alpha j}\in\C$ for $0\leq j<\mult_\alpha$, and blocks
    $A_\alpha$ of bond dimension $\dim_\alpha$. The vertical direct-sum tensor is the
    block-diagonal direct sum over the flattened index $(\alpha,j)$ that repeats
    $A_\alpha$ with weight $\omega_{\alpha j}$, written
    $\bigoplus_\alpha\bigoplus_j\omega_{\alpha j}A_\alpha$. Its matrix product
    vector decomposes as
    \begin{align}
        V^{(N)}\Bigl(\textstyle\bigoplus_\alpha\bigoplus_j
        \omega_{\alpha j}A_\alpha\Bigr)(\sigma)
        &=\sum_\alpha\sum_j(\omega_{\alpha j})^N
          V^{(N)}(A_\alpha)(\sigma).
        \label{eq:mpdo_vertical_mpv_sum}
    \end{align}
\end{lemma}
\begin{proof}\leanok
    \uses{thm:mpv_decomposition}
    The vertical direct-sum tensor is the block-diagonal direct sum over the single
    flattened index $q$, with flattened weights $\widetilde\omega_q$ and blocks
    $\widetilde A_q$. By the block matrix-product-vector decomposition
    (Theorem~\ref{thm:mpv_decomposition}),
    \begin{align}
        V^{(N)}\Bigl(\textstyle\bigoplus_q
        \widetilde\omega_q\widetilde A_q\Bigr)(\sigma)
        &=\sum_q(\widetilde\omega_q)^N
          V^{(N)}(\widetilde A_q)(\sigma).
        \notag
    \end{align}
    The flattening identifies $q$ with the pair $(\alpha,j)$ via the bijection
    \begin{align}
        \{(\alpha,j):0\leq\alpha<g,\ 0\leq j<\mult_\alpha\}
        &\simeq
        \{0,\ldots,\textstyle\sum_\alpha\mult_\alpha-1\},
        \notag
    \end{align}
    under which $\widetilde\omega_q=\omega_{\alpha j}$ and
    $\widetilde A_q=A_\alpha$. Reindexing the sum along this bijection yields the
    double sum over $(\alpha,j)$.
\end{proof}

\begin{lemma}[Regrouping a block family by repeated representatives]
    \label{lem:vertical_grouping_equiv}
    \lean{MPOTensor.sameMPV₂Pos_toTensorFromBlocks_verticalAssembledTensor_of_equiv}
    \leanok
    \uses{lem:mpv_vertical_assembled, thm:mpv_decomposition}
    Let $(B_k)_{k\in I}$ be a finite family of blocks with weights $\nu_k$,
    and let $(A_\alpha)_{\alpha\in J}$ be a finite family of representative
    blocks. Suppose there is a bijection
    \begin{align}
        e:I&\simeq\{(\alpha,j):\alpha\in J,\ 0\leq j<m_\alpha\}
        \notag
    \end{align}
    such that, writing $e(k)=(\alpha,j)$, one has
    $\nu_k=\omega_{\alpha j}$ and, for every positive length $N$,
    $V^{(N)}(B_k)=V^{(N)}(A_\alpha)$. Then
    \begin{align}
        V^{(N)}\left(\textstyle\bigoplus_{k\in I}\nu_kB_k\right)
        &=
        V^{(N)}\Bigl(      \textstyle\bigoplus_{\alpha\in J}
          \bigoplus_{j=0}^{m_\alpha-1}\omega_{\alpha j}A_\alpha
        \Bigr)
        \notag
    \end{align}
    for every positive length $N$. This is the finite-sum
    reindexing after the vertical sectors have been grouped in
    \cite[Proposition~4.13]{Cirac2016MPDO_arXiv}.
\end{lemma}
\begin{proof}\leanok
    \uses{lem:mpv_vertical_assembled, thm:mpv_decomposition}
    For $N>0$ and every configuration $\sigma$, the two sides expand as
    \begin{align}
        \sum_{k\in I}\nu_k^NV^{(N)}(B_k)(\sigma),
        \notag\\
        \sum_{\alpha\in J}\sum_{j=0}^{m_\alpha-1}
        \omega_{\alpha j}^NV^{(N)}(A_\alpha)(\sigma).
        \notag
    \end{align}
    The summand indexed by $k$ in the first sum equals the summand indexed by
    $e(k)$ in the second, so reindexing by $e$ proves the equality.
\end{proof}

\begin{theorem}[Invariant vertical projections in literal CPSV canonical form]
    \label{thm:mpdo_vertical_projector_closure_literal_cpsv}
    \lean{MPSTensor.CPSVCanonicalFormData.ActiveBNTRefinement.%
        hasInvariantProjectorClosure_verticalTensor}
    \lean{MPSTensor.IsCPSVCanonicalForm.hasInvariantProjectorClosure_verticalTensor}
    \leanok
    \uses{def:mpdo, def:cpsv_canonical_form, def:vertical_tensor_mpdo,
        def:invariant_projector_closure,
        thm:representative_one_sub_mp_zero,
        thm:cpsv_original_coordinate_marked_lemma_l}
    Let $M$ generate matrix product density operators, and suppose that its
    doubled-index MPS tensor is in literal CPSV canonical form. If an
    orthogonal projection $P$ satisfies
    $P\widetilde M_v=P\widetilde M_vP$ for every vertical letter $v$, then
    \begin{align}
        \widetilde M_vP&=P\widetilde M_vP
        \label{eq:cpsv_literal_vertical_projector_closure}
    \end{align}
    for every $v$. Thus the vertically viewed tensor has
    invariant-projector closure.
\end{theorem}

\begin{proof}\leanok
    \uses{thm:representative_one_sub_mp_zero,
        thm:cpsv_original_coordinate_marked_lemma_l}
    The left-invariance hypothesis gives
    $P_1H^{(N)}=P_1H^{(N)}P_1$ at every positive length. Since $H^{(N)}$ is
    positive and $P=P^\dagger$,
    \begin{align}
        P_1H^{(N)}&=P_1H^{(N)}P_1=H^{(N)}P_1.
        \label{eq:cpsv_literal_first_site_commutation}
    \end{align}
    The matrix coefficients of
    (\ref{eq:cpsv_literal_first_site_commutation}) give agreement of the
    corresponding first-site actions. The original-coordinate marked
    Lemma~L gives equality of their inserted tensors on the original bond
    space. Reading this equality as a vertical one-site action yields
    (\ref{eq:cpsv_literal_vertical_projector_closure}).
\end{proof}

\begin{theorem}[Invariant vertical projections are reducing]
    \label{thm:mpdo_vertical_projector_closure}
    \lean{MPOTensor.IsHorizontalCF.hasInvariantProjectorClosure_verticalTensor}
    \leanok
    \uses{def:mpdo, def:horizontal_cf_mpdo, def:vertical_tensor_mpdo,
        def:invariant_projector_closure,
        thm:representative_one_sub_mp_zero}
    Let $M$ be in normalized BNT-refined horizontal form and generate matrix
    product density operators. If an orthogonal projection $P$ satisfies
    $P\widetilde M_v=P\widetilde M_vP$ for every vertical letter $v$, then it
    also satisfies $\widetilde M_vP=P\widetilde M_vP$ for every vertical
    letter $v$. Thus the vertically viewed tensor has invariant-projector
    closure.
\end{theorem}

\begin{proof}
    \leanok
    \uses{thm:representative_one_sub_mp_zero}
    Translate the left-invariance identity into the one-site ket and bra
    actions of the original tensor. The normalized BNT-refined version of
    Lemma~L gives the opposite one-sided identity. Translating back to the
    vertically viewed tensor proves the claim.
\end{proof}

\begin{theorem}[Irreducible reducing sectors of the vertically viewed tensor]
    \label{thm:mpdo_vertical_irreducible_reducing_sectors}
    \lean{MPOTensor.IsHorizontalCF.exists_irreducible_verticalBlockDecomp_with_isometry}
    \leanok
    \uses{thm:mpdo_vertical_projector_closure,
        thm:projector_closure_irreducible_corners}
    Under the hypotheses of
    Theorem~\ref{thm:mpdo_vertical_projector_closure}, there are positive
    integers $d_k$, tensors $B_k$ of bond dimension $d_k$, and isometries
    $V_k:\C^{d_k}\to\C^d$ such that
    \begin{align}
        V_k^\dagger V_k&=\Id,
        \notag\\
        \sum_kV_kV_k^\dagger&=\Id,
        \notag\\
        V_k^\dagger V_l&=0\quad(k\ne l).
        \notag
    \end{align}
    Every $B_k$ is irreducible, and for every vertical letter $v$,
    \begin{align}
        \widetilde M_vV_k&=V_kB_k^v,
        \notag\\
        V_k^\dagger\widetilde M_v&=B_k^vV_k^\dagger,
        \notag\\
        B_k^v&=V_k^\dagger\widetilde M_vV_k.
        \notag
    \end{align}
    Moreover, $\widetilde M$ and the unit-weight direct sum
    $\bigoplus_kB_k$ have the same matrix product vectors at every positive
    length.
\end{theorem}

\begin{proof}
    \leanok
    \uses{thm:mpdo_vertical_projector_closure,
        thm:projector_closure_irreducible_corners}
    Apply the irreducible-corner decomposition under invariant-projector
    closure to $\widetilde M$.
\end{proof}

\begin{theorem}[Complete vertical reducing projectors]
    \label{thm:mpdo_vertical_complete_reducing_projectors}
    \lean{MPOTensor.IsHorizontalCF.exists_complete_verticalReducingProjectors}
    \leanok
    \uses{thm:mpdo_vertical_irreducible_reducing_sectors}
    Under the same hypotheses, there are orthogonal projections $P_k$ that
    are complete and pairwise orthogonal and, for every vertical letter $v$,
    reduce the corresponding letter:
    \begin{align}
        \sum_kP_k&=\Id,
        \notag\\
        P_kP_l&=0\quad(k\ne l),
        \notag\\
        \widetilde M_vP_k&=P_k\widetilde M_v.
        \notag
    \end{align}
\end{theorem}

\begin{proof}
    \leanok
    \uses{thm:mpdo_vertical_irreducible_reducing_sectors}
    For the isometries of
    Theorem~\ref{thm:mpdo_vertical_irreducible_reducing_sectors}, set
    $P_k=V_kV_k^\dagger$. The isometry, orthogonality, completeness, and
    intertwining identities give the four assertions directly.
\end{proof}

\begin{figure}[ht]
    \centering
    $U\mathcal K U^\dagger=\bigoplus_\alpha M_\alpha$.
    \caption{The vertically viewed tensor decomposes through complete
    orthogonal physical sectors. Each $V_k$ embeds the irreducible tensor
    $B_k$, while $P_k=V_kV_k^\dagger$ is its reducing range projection. The
    displayed identity is equivalent, for every vertical letter $v$, to
    $\widetilde M_v=\sum_kV_kB_{k,v}V_k^\dagger$.
    This is the iterated invariant-subspace construction in
    \cite[lines~200--219]{Cirac2016MPDO_arXiv}, applied to the reduction in
    \cite[Proposition~4.13, lines~1873--1887]{Cirac2016MPDO_arXiv}.}
    \label{fig:mpdo_vertical_reducing_sectors}
\end{figure}

\begin{theorem}[Spectrally normalized vertical sectors with physical isometries]
    \label{thm:mpdo_vertical_normal_sectors_with_isometries}
    \lean{MPOTensor.IsHorizontalCF.exists_normal_verticalBlockDecomp_with_isometry}
    \leanok
    \uses{thm:mpdo_vertical_irreducible_reducing_sectors,
        thm:mpdo_vertical_no_periodic_vectors_horizontal_cf,
        lem:irreducible_transfer_perron_pair,
        lem:irreducible_block_spectral_normalization}
    Under the same hypotheses, there are positive integers $d_k$, positive
    scalars $\nu_k$, normal tensors $B_k$ of bond dimension $d_k$, and
    isometries $V_k:\C^{d_k}\to\C^d$ with pairwise orthogonal ranges. For
    every sector $k$ and every vertical letter $v$,
    \begin{align}
        \widetilde M_vV_k&=V_k(\nu_kB_k^v),
        \notag\\
        V_k^\dagger\widetilde M_v&=(\nu_kB_k^v)V_k^\dagger,
        \notag\\
        \nu_kB_k^v&=V_k^\dagger\widetilde M_vV_k.
        \notag
    \end{align}
    Moreover, every vertical letter is reconstructed literally by
    \begin{align}
        \widetilde M_v
        &=\sum_kV_k(\nu_kB_k^v)V_k^\dagger.
        \label{eq:mpdo_vertical_normal_reconstruction}
    \end{align}
\end{theorem}

\begin{proof}
    \leanok
    \uses{thm:mpdo_vertical_irreducible_reducing_sectors,
        thm:mpdo_vertical_no_periodic_vectors_horizontal_cf,
        lem:irreducible_transfer_perron_pair,
        lem:irreducible_block_spectral_normalization}
    Retain exactly the irreducible corners containing a nonzero letter. The
    Perron--Frobenius eigenvalue of each retained corner is positive, and the
    absence of nontrivial periodic vectors makes the inverse-square-root
    rescaling normal. Put the removed square root into $\nu_k$ and reindex the
    retained corners. The original complete reducing decomposition expresses
    each vertical letter as the sum of all compressed corners; every omitted
    summand vanishes, which gives
    (\ref{eq:mpdo_vertical_normal_reconstruction}). Thus the retained range
    projections need not sum to the identity, but their orthogonal complement
    is a zero sector.
\end{proof}

\begin{theorem}[Vertical normal sectors grouped into BNT representatives]
    \label{thm:mpdo_vertical_bnt_grouping_with_isometries}
    \lean{MPOTensor.IsHorizontalCF.exists_verticalBNTGrouping_with_isometry}
    \leanok
    \uses{def:mpdo, def:horizontal_cf_mpdo, def:vertical_tensor_mpdo,
        def:bnt_cpsv, def:mpv_phase_class_data}
    Under the same hypotheses, partition the normalized vertical sectors by
    equality of their matrix product vector families up to a fixed nonzero
    scalar per site. There are normal representatives $M_\alpha$, an enumeration
    $k=(\alpha,q)$ of the original sectors, invertible matrices
    $X_{\alpha,q}$, and complex numbers $\zeta_{\alpha,q}$ of unit modulus,
    with $X_{\alpha,1}=\Id$ and $\zeta_{\alpha,1}=1$, such that
    \begin{align}
        B_{\alpha,q}^v
        &=\zeta_{\alpha,q}X_{\alpha,q}
          M_\alpha^vX_{\alpha,q}^{-1}.
        \label{eq:mpdo_vertical_gauge_phase}
    \end{align}
    The representatives are pairwise gauge-phase distinct and form a basis of
    normal tensors for the weighted block tensor
    $\bigoplus_k\nu_kB_k$. Every grouped coefficient
    $\nu_{\alpha,q}\zeta_{\alpha,q}$ is positive. With the physical isometries
    from Theorem~\ref{thm:mpdo_vertical_normal_sectors_with_isometries}
    retained unchanged, one has both reducing identities, exact compression,
    and the literal reconstruction
    \begin{align}
        \widetilde M_v
        &=\sum_\alpha\sum_q
          V_{\alpha,q}
          \bigl(\nu_{\alpha,q}\zeta_{\alpha,q}
          X_{\alpha,q}M_\alpha^vX_{\alpha,q}^{-1}\bigr)
          V_{\alpha,q}^\dagger.
        \label{eq:mpdo_vertical_grouped_reconstruction}
    \end{align}
    For each pair $(\alpha,q)$, the range projector of $V_{\alpha,q}$ obeys
    the representative-loop identity and gives a nonzero compression at some
    finite length. This is the decomposition and positivity argument of
    \cite[Proposition~4.13, lines 1898--1902]{Cirac2016MPDO_arXiv}. The
    relative Gram identity and proportional-unitary normalization of
    lines~1903--1908 are not asserted here.
\end{theorem}

\begin{proof}\leanok
    \uses{thm:mpdo_vertical_normal_sectors_with_isometries,
        thm:cpsv_normal_mpv_phase_gauge,
        thm:cpsv_normal_separated_eventually_li,
        thm:mpdo_grouped_sector_compression_nonzero,
        def:mpdo_sector_projector, thm:mpdo_sector_weight_pos}
    Take one representative from each phase class.
    Theorem~\ref{thm:cpsv_normal_mpv_phase_gauge} gives equality of bond
    dimensions and (\ref{eq:mpdo_vertical_gauge_phase}) for every copy.
    Distinct representatives are gauge-phase separated by construction, so
    Theorem~\ref{thm:cpsv_normal_separated_eventually_li} gives eventual linear
    independence. The phase-class regrouping of the finite sum gives the BNT
    expansion. The representative copy is assigned the identity gauge and
    phase one, while the self-overlap normalization of normal tensors shows
    that every phase has unit modulus. Substitution into the two intertwining
    identities, the compression identity, and the reconstruction
    (\ref{eq:mpdo_vertical_normal_reconstruction}) proves all remaining
    algebraic assertions without altering any physical isometry. Cyclicity of
    trace gives the sector-loop identities. The nonzero-corner separation
    theorem supplies a nonzero compression for every sector, and the sector
    trace identity and MPDO positivity then make every grouped coefficient
    positive.
\end{proof}

\begin{definition}[Block-row map of sector isometries]
    \label{def:vertical_sector_block_row}
    \lean{Matrix.sigmaBlockRow}
    \leanok
    Let $W_k:\C^{D_k}\to\C^d$ be a finite family of rectangular
    matrices. Their block-row map
    $U:\C^d\to\bigoplus_k\C^{D_k}$ is defined by
    $Ux=(W_k^\dagger x)_k$. This realizes the coisometry in the conclusion of
    \cite[Proposition~4.13, lines~1863--1870]{Cirac2016MPDO_arXiv} once the
    sector isometries have been normalized.
\end{definition}

\begin{theorem}[Conjugation and reconstruction from orthogonal sectors]
    \label{thm:vertical_sector_coisometry}
    \lean{Matrix.sigmaBlockRow_isCoisometry}
    \lean{Matrix.sigmaBlockRow_conjugation}
    \lean{Matrix.sigmaBlockRow_reconstruction}
    \leanok
    \uses{def:vertical_sector_block_row}
    Suppose that
    \begin{align}
        W_k^\dagger W_k&=\Id,
        \notag\\
        W_k^\dagger W_l&=0\quad(k\ne l),
        \notag
    \end{align}
    and, for every letter $v$ and every sector $k$,
    $T_vW_k=W_kB_k^v$. Then the block-row map $U$ is a coisometry and
    conjugates $T$ to the direct sum of the sector tensors:
    \begin{align}
        UU^\dagger&=\Id,
        \notag\\
        UT_vU^\dagger&=\bigoplus_kB_k^v.
        \notag
    \end{align}
    If the letters also satisfy the literal sector decomposition
    \begin{align}
        T_v&=\sum_kW_kB_k^vW_k^\dagger,
        \label{eq:mpdo_sector_literal_decomposition}
    \end{align}
    then
    \begin{align}
        T_v&=U^\dagger\!\left(\bigoplus_kB_k^v\right)U.
        \notag
    \end{align}
    This last equality is a conditional consequence of
    (\ref{eq:mpdo_sector_literal_decomposition}); no additional assertion
    from the source is used.
\end{theorem}

\begin{proof}
    \leanok
    The $(k,l)$ block of $UU^\dagger$ is $W_k^\dagger W_l$, hence the
    first identity follows from orthonormality. The corresponding block of
    $UT_vU^\dagger$ is
    $W_k^\dagger T_vW_l=W_k^\dagger W_lB_l^v$, which is $B_k^v$ for $k=l$
    and zero otherwise. Finally, multiplication of the block diagonal matrix
    by the block row gives the sum $\sum_kW_kB_k^vW_k^\dagger$, which is
    $T_v$ by hypothesis. The row indices of $U$ form the dependent disjoint
    union $\coprod_k\{0,\ldots,D_k-1\}$. Its canonical order-preserving
    bijection with $\{0,\ldots,\sum_kD_k-1\}$ gives the coordinate reindexing
    required by the vertical canonical-form definition.
    % Lean adapter: this final reindexing is `finSigmaFinEquiv`.
\end{proof}

\begin{theorem}[Vertical canonical form from orthogonal grouped sectors]
    \label{thm:vertical_cf_of_grouped_orthogonal_sectors}
    \lean{MPOTensor.isVerticalCF_of_grouped_orthogonal_sectors}
    \leanok
    \uses{def:vertical_tensor_mpdo, def:vertical_cf_mpdo, def:bnt}
    Let $\{A_\alpha\}_{\alpha=0}^{g-1}$ be a BNT for the vertically viewed
    tensor $\widetilde M$. For each $\alpha$, let $r_\alpha>0$, let
    $\omega_{\alpha,q}>0$ for $0\leq q<r_\alpha$, and let
    $W_{\alpha,q}:\C^{D_\alpha}\to\C^d$ satisfy
    \begin{align}
        W_{\alpha,q}^\dagger W_{\alpha,q}&=\Id,
        \notag\\
        W_{\alpha,q}^\dagger W_{\beta,p}&=0
        \quad((\alpha,q)\ne(\beta,p)).
        \notag
    \end{align}
    Suppose that, for every vertical letter $v$ and every pair $(\alpha,q)$,
    $\widetilde M_vW_{\alpha,q}
    =W_{\alpha,q}(\omega_{\alpha,q}A_\alpha^v)$, and that the letters have
    the exact sector decomposition
    \begin{align}
        \widetilde M_v
        &=\sum_{\alpha=0}^{g-1}\sum_{q=0}^{r_\alpha-1}
          W_{\alpha,q}(\omega_{\alpha,q}A_\alpha^v)
          W_{\alpha,q}^\dagger.
        \notag
    \end{align}
    Then $M$ is in vertical canonical form.
    This is the final finite-dimensional implication in
    \cite[Proposition~4.13, lines~1863--1870]{Cirac2016MPDO_arXiv}.
\end{theorem}

\begin{proof}
    \leanok
    \uses{thm:vertical_sector_coisometry}
    Form the block-row map $U$ from the adjoints
    $W_{\alpha,q}^\dagger$. Theorem~\ref{thm:vertical_sector_coisometry}
    gives
    \begin{align}
        UU^\dagger&=\Id,
        \notag\\
        U\widetilde M_vU^\dagger
        &=\bigoplus_{\alpha,q}\omega_{\alpha,q}A_\alpha^v,
        \notag\\
        \widetilde M_v
        &=U^\dagger\!\left(\bigoplus_{\alpha,q}\omega_{\alpha,q}A_\alpha^v\right)U.
        \notag
    \end{align}
    Identifying the dependent pairs $(\alpha,q)$ and their bond indices with
    the standard coordinates of the two finite direct sums gives the required
    matrix dimensions. Positivity of the multiplicities and weights, together
    with the BNT hypothesis, supplies the remaining conditions of vertical
    canonical form.
\end{proof}

\begin{lemma}[A retained vertical phase class exists]
    \label{lem:mpdo_vertical_retained_class_nonempty}
    \lean{MPOTensor.IsHorizontalCF.exists_nonempty_verticalBNTGrouping}
    \leanok
    \uses{def:horizontal_cf_mpdo,
        thm:mpdo_vertical_bnt_grouping_with_isometries,
        lem:mpv_phase_class_nonempty}
    Under the hypotheses of
    Theorem~\ref{thm:mpdo_vertical_bnt_grouping_with_isometries}, the retained
    vertical decomposition has at least one sector and hence at least one
    phase class.
\end{lemma}

\begin{proof}
    \leanok
    The unit-modulus weight witness in the normalized BNT-refined horizontal
    form, together with left canonicity of its representative, makes the doubled-index tensor
    nonzero. Hence its vertical reshaping $\widetilde M$ is nonzero. If the
    retained decomposition had no sector,
    (\ref{eq:mpdo_vertical_grouped_reconstruction}) would make every letter of
    $\widetilde M$ zero, a contradiction. Passing from the nonempty finite set
    of retained sectors to its phase partition gives a retained representative.
\end{proof}

\begin{theorem}[Nonvanishing of the vertical reshaping]
    \label{thm:mpdo_vertical_tensor_nonzero}
    \lean{MPOTensor.IsHorizontalCF.verticalTensor_ne_zero}
    \leanok
    \uses{def:horizontal_cf_mpdo, def:vertical_tensor_mpdo}
    If an MPO tensor $M$ is in normalized BNT-refined horizontal form, then
    its vertical reshaping is nonzero: $\widetilde M\ne0$.
    This is the nonvanishing observation used in
    \cite[Proposition~4.13, lines~1895--1902]{Cirac2016MPDO_arXiv}.
\end{theorem}

\begin{proof}
    \leanok
    \uses{def:horizontal_cf_mpdo, def:vertical_tensor_mpdo,
        lem:sector_bnt_weight_unit_exists_of_struct}
    Write the doubled-index tensor in normalized BNT-refined horizontal form
    as $\widehat M=GSG^{-1}$, where $S$ is the weighted direct sum of its basis
    tensors. Choose a copy whose weight has unit modulus. Its basis tensor
    cannot have every letter equal to zero, since left canonicity gives
    $\sum_i(A^i)^\dagger A^i=\Id$.
    The corresponding nonzero weight and nonzero letter give $S\ne0$, hence
    $\widehat M\ne0$. If $\widetilde M=0$, the bijection between its vertical
    index and the two horizontal bond indices would give $\widehat M=0$, a
    contradiction.
\end{proof}

\begin{theorem}[Algebraic BNT from a grouped vertical decomposition]
    \label{thm:mpdo_vertical_grouped_is_bnt}
    \lean{MPOTensor.isBNT_verticalTensor_of_grouping}
    \leanok
    \uses{def:vertical_tensor_mpdo, def:bnt,
        def:bnt_cpsv, def:mpv_phase_class_data}
    Let $M$ be an MPO tensor and put $T=\widetilde M$. Let $B_k$ have bond
    dimension $d_k$, let $\mu_k\in\C$, and let
    $V_k:\C^{d_k}\to\C^d$ be an isometry for $0\leq k<r$. Assume, for every
    vertical letter $v$, that
    \begin{align}
        V_k^\dagger T_v&=(\mu_kB_k^v)V_k^\dagger,
        \notag\\
        T_v&=\sum_{k=0}^{r-1}V_k(\mu_kB_k^v)V_k^\dagger.
        \notag
    \end{align}
    Partition the tensors $B_k$ into matrix-product-vector phase classes and
    choose one representative $A_\alpha$ from each class. If
    $\{A_\alpha\}_\alpha$ is a basis of normal tensors in the spectral sense
    for the weighted direct sum
    $A_{\mathrm{ret}}=\bigoplus_{k=0}^{r-1}\mu_kB_k$, then
    $\{A_\alpha\}_\alpha$ is an algebraic basis of normal tensors for $T$.

    This is the basis-of-normal-tensors conclusion in
    \cite[Proposition~4.13, lines~1863--1921]{Cirac2016MPDO_arXiv}, stated for
    the grouped decomposition before the final coisometry is formed.
\end{theorem}

\begin{proof}
    \leanok
    \uses{thm:cpsv_bnt_is_bnt, lem:bnt_of_same_mpv2_pos,
        thm:mpv_decomposition}
    Put $C_k^v=\mu_kB_k^v$. Induction on a word $w$ gives
    $V_k^\dagger T^w=C_k^wV_k^\dagger$.
    For a nonempty word $w=(v,w')$, reconstruction of its first letter,
    cyclicity of trace, and $V_k^\dagger V_k=\Id$ therefore give
    \begin{align}
        \tr(T^w)
        &=\sum_k\tr\!\left(V_kC_k^vV_k^\dagger T^{w'}\right)
        \notag\\
        &=\sum_k\tr(C_k^vC_k^{w'})
          =\sum_k\mu_k^{|w|}\tr(B_k^w)
          =V^{(|w|)}(A_{\mathrm{ret}})_w.
        \notag
    \end{align}
    Thus $T$ and $A_{\mathrm{ret}}$ have the same matrix product vectors at
    every positive length. The spectral BNT hypothesis gives an algebraic BNT
    for $A_{\mathrm{ret}}$ by Theorem~\ref{thm:cpsv_bnt_is_bnt}; in particular,
    it supplies the required positive bond dimensions.
    Lemma~\ref{lem:bnt_of_same_mpv2_pos} then gives the same algebraic BNT for
    $T$.
\end{proof}

\begin{theorem}[Pairwise Figure~8 identity under literal canonical form]
    \label{thm:cpsv_literal_pairwise_vertical_gram_conjugation}
    \lean{MPSTensor.IsCPSVCanonicalForm.gramDressing_eq_of_two_grouped_corners}
    \leanok
    \uses{def:mpdo, def:cpsv_canonical_form,
        thm:mpdo_reflected_marked_chain,
        thm:cpsv_original_coordinate_marked_lemma_l}
    Let $M$ be an MPO tensor whose doubled-index MPS tensor is in the literal
    CPSV canonical form, and suppose that $M$ generates positive semidefinite
    operators on every nonempty chain. Let $A$ be a tensor, let $X$ and $Y$ be
    invertible matrices, and suppose that two vertical corners satisfy
    \begin{align}
        V_X^\dagger\widetilde M^vV_X
        &=c_X X A^vX^{-1},
        &
        c_X&>0,
        \notag\\
        V_Y^\dagger\widetilde M^vV_Y
        &=c_Y Y A^vY^{-1},
        &
        c_Y&>0.
        \notag
    \end{align}
    Then the Gram-dressed tensors agree:
    \begin{align}
        (X^\dagger X)A^v(X^\dagger X)^{-1}
        &=(Y^\dagger Y)A^v(Y^\dagger Y)^{-1}
        \label{eq:cpsv_literal_pairwise_gram_conj}
    \end{align}
    for every vertical letter $v$.
\end{theorem}

\begin{proof}\leanok
    \uses{thm:mpdo_reflected_marked_chain,
        thm:cpsv_original_coordinate_marked_lemma_l}
    For $Z=X,Y$, set
    \begin{align}
        f^Z_{(r,s),(i,j)}
        &=c_Z^{-1}\overline{(V_Z Z)_{ir}}
          \bigl(V_Z(Z^{-1})^\dagger\bigr)_{js}.
        \notag
    \end{align}
    The corner identity identifies the mark with coefficients $f^Z$ with the
    horizontal slice of $(Z^\dagger Z)A^v(Z^\dagger Z)^{-1}$.
    Self-adjointness of the two corner compressions identifies both closed
    marked chains with the same reflected chain. The original-coordinate
    marked Lemma~L for literal canonical form therefore gives equality of the
    two marks, hence (\ref{eq:cpsv_literal_pairwise_gram_conj}).
\end{proof}

\begin{theorem}[Positive relative Gram scalar under literal canonical form]
    \label{thm:cpsv_literal_positive_relative_gram_scalar}
    \lean{MPSTensor.IsCPSVCanonicalForm.gram_eq_pos_smul_gram_of_two_grouped_corners}
    \leanok
    \uses{thm:cpsv_literal_pairwise_vertical_gram_conjugation,
        thm:normal_tensor_gram_rigidity}
    Let $M$ be an MPO tensor whose doubled-index MPS tensor is in the literal
    CPSV canonical form, and suppose that $M$ generates positive semidefinite
    operators on every nonempty chain. Let $A$ be normal, let $X$ and $Y$ be
    invertible matrices, and suppose that
    \begin{align}
        V_X^\dagger\widetilde M^vV_X
        &=c_X X A^vX^{-1},
        &
        c_X&>0,
        \notag\\
        V_Y^\dagger\widetilde M^vV_Y
        &=c_Y Y A^vY^{-1},
        &
        c_Y&>0.
        \notag
    \end{align}
    Then there is a real number $\omega>0$ such that
    \begin{align}
        X^\dagger X&=\omega\,Y^\dagger Y.
        \label{eq:cpsv_literal_positive_relative_gram}
    \end{align}
\end{theorem}

\begin{proof}\leanok
    \uses{thm:cpsv_literal_pairwise_vertical_gram_conjugation,
        thm:normal_tensor_gram_rigidity}
    Equation~(\ref{eq:cpsv_literal_pairwise_gram_conj}) makes the ratio of the
    two Gram matrices commute with every letter of $A$. Normality makes this
    ratio scalar, and positive definiteness makes the scalar a positive real
    number.
\end{proof}

% These two theorems concern Figures 7--8 in the proof of Proposition 4.13.
% They are not the Appendix C.4 positive-fusion statement tracked by #4648.
% Source anchor: CPSV16 Proposition 4.13, statement and proof at lines 1863--1922.
% Local fix (rectangular coisometry and reconstruction): the source calls U an
% isometry and prints only the compressed identity. See
% docs/paper-gaps/cpgsv17_vertical_isometry_zero_sector.tex.
\begin{theorem}[Vertical canonical form of a canonical MPDO]
    \label{thm:vertical_cf_of_cpsv_canonical_form}
    \notready
    \uses{def:mpdo, def:cpsv_canonical_form, def:vertical_cf_mpdo}
    Let $M$ be an MPO tensor whose doubled-index MPS tensor is in the literal
    CPSV canonical form. If $M$ generates matrix product density operators,
    then $M$ is in canonical form in the vertical direction. More precisely,
    there is a basis of normal tensors $\{M_\alpha\}_\alpha$ for the vertically
    viewed tensor $\widetilde M$, together with positive diagonal matrices
    $\mu_\alpha$ and a coisometry $U$ from the physical space onto the retained
    nonzero sectors. Writing
    $B=\bigoplus_\alpha\mu_\alpha\otimes M_\alpha$, one has
    \begin{align}
        UU^\dagger&=\Id.
        \label{eq:cpsv_prop_4_13_coisometry}
    \end{align}
    \begin{align}
        U\widetilde M U^\dagger&=B.
        \label{eq:cpsv_prop_4_13_vertical_form}
    \end{align}
    \begin{align}
        \widetilde M&=U^\dagger BU.
        \label{eq:cpsv_prop_4_13_reconstruction}
    \end{align}
    The source explicitly displays
    (\ref{eq:cpsv_prop_4_13_vertical_form}) and calls $U$ an isometry
    \cite[Proposition~4.13]{Cirac2016MPDO_arXiv}. In this rectangular
    orientation, however, $U$ maps the physical space onto the retained sector
    space and satisfies (\ref{eq:cpsv_prop_4_13_coisometry}); thus $U$ is a
    coisometry and $U^\dagger$ is the isometric inclusion. The reverse identity
    (\ref{eq:cpsv_prop_4_13_reconstruction}) is not displayed separately in
    the proposition. It records the exact reconstruction implicit in the
    preceding assertion that $\widetilde M$ is itself in canonical form,
    including the permitted zero orthogonal complement.
\end{theorem}

\begin{proof}[Proof outline]
    \uses{thm:cpsv_canonical_form_active_bnt_refinement,
        thm:cpsv_active_representative_sector_decomposition,
        thm:cpsv_full_coordinate_marked_trace_identity,
        thm:cpsv_retained_representative_grouped_lemma_l,
        thm:cpsv_original_coordinate_marked_lemma_l,
        thm:mpdo_vertical_projector_closure_literal_cpsv,
        thm:projector_closure_irreducible_corners,
        thm:mpdo_no_periodic_sector_projector,
        thm:psd_commute_of_commute_pow,
        lem:irreducible_transfer_perron_pair,
        lem:irreducible_block_spectral_normalization,
        thm:cpsv_normal_mpv_phase_gauge,
        thm:cpsv_normal_separated_eventually_li,
        thm:mpdo_sector_compression_trace_pos,
        thm:mpdo_sector_weight_pos,
        thm:cpsv_literal_pairwise_vertical_gram_conjugation,
        thm:cpsv_literal_positive_relative_gram_scalar,
        thm:normal_tensor_gram_rigidity,
        thm:vertical_sector_coisometry}
    Write $A$ for the doubled-index MPS tensor of $M$. If $A=0$, take no
    retained sectors; the unique map onto the zero-dimensional sector space
    gives the conclusion. Suppose henceforth that $A\ne0$. Apply
    Theorem~\ref{thm:cpsv_canonical_form_active_bnt_refinement} to its literal
    canonical-form data. The active listed indices are identified with pairs
    $(\alpha,k)$ consisting of a phase class and a copy, and every active copy
    retains its original nonzero coefficient $\nu_{\alpha,k}$, its bond
    dimension, a unit phase, and an invertible gauge $X_{\alpha,k}$. The
    chosen tensors $C_\alpha$ form a basis of normal tensors. Adjoin the
    inactive zero-weight indices to these pairs and call the resulting grouped
    index set $\mathcal G$. There is an explicit equivalence from the dependent
    coordinates of $\mathcal G$ to the original flattened listed coordinates.
    If $T_{\mathrm{grp}}$ is the grouped direct sum, reindexed along this
    equivalence, then there are a block-diagonal gauge $G$ and the original
    coisometry $V$ such that
    \begin{align}
        \bigoplus_\ell\mu_\ell A_\ell^i
        &=GT_{\mathrm{grp}}^iG^{-1},
        \label{eq:cpsv_prop_4_13_active_bnt}
    \end{align}
    and
    \begin{align}
        A^i&=V^\dagger GT_{\mathrm{grp}}^iG^{-1}V.
        \label{eq:cpsv_prop_4_13_active_reconstruction}
    \end{align}
    \begin{align}
        VV^\dagger&=\Id.
        \label{eq:cpsv_prop_4_13_original_coisometry}
    \end{align}
    The active summands of $T_{\mathrm{grp}}$ are
    $\nu_{\alpha,k}\zeta_{\alpha,k}C_\alpha^i$, with
    $|\zeta_{\alpha,k}|=1$. Every inactive summand is the original listed block
    with coefficient zero. Thus the grouped coordinates are exact without
    deleting the inactive complement or introducing a selecting coisometry.

    Let $H^{(N)}$ be the positive operator generated horizontally by
    $\widetilde M$. If
    an orthogonal projector $P$ on the physical space satisfies
    $P\widetilde M=P\widetilde MP$, and $P_1$ acts on the first site, then
    self-adjointness gives
    \begin{align}
        P_1H^{(N)}
        &=P_1H^{(N)}P_1
          =(P_1H^{(N)}P_1)^\dagger
          =H^{(N)}P_1.
        \label{eq:cpsv_prop_4_13_first_site}
    \end{align}
    Write $Y_R(P)$ and $Y_C(P)$ for the right and central first-site
    contractions determined by $P$. The active representative sector
    decomposition has weights
    $\lambda_{\alpha,k}=\nu_{\alpha,k}\zeta_{\alpha,k}$, all nonzero, and its
    representatives have simultaneous word-tuple separation at every
    sufficiently large length. It has the same matrix product vectors as
    $T_{\mathrm{grp}}$ at every positive length. Hence
    Theorem~\ref{thm:cpsv_retained_representative_grouped_lemma_l} gives
    \begin{align}
        Y_R(P)C_\alpha^i&=Y_C(P)C_\alpha^i
        \label{eq:cpsv_prop_4_13_active_separation}
    \end{align}
    on every active representative. Theorem~\ref{thm:cpsv_original_coordinate_marked_lemma_l}
    then uses the listed block gauge and the ambient coisometry, with
    $VV^\dagger=\Id$, to give the original-space inserted-tensor equality.
    Theorem~\ref{thm:mpdo_vertical_projector_closure_literal_cpsv} therefore
    yields
    \begin{align}
        \widetilde MP&=P\widetilde MP.
        \label{eq:cpsv_prop_4_13_projector_closure}
    \end{align}
    Thus every one-sided invariant subspace is reducing. The inactive listed
    coordinates remain present with weight zero throughout this argument; no
    insertion identity is asserted for an individual inactive block.

    % Local fix (periodic-sector quantifier): use one noncommuting chain length,
    % not the all-length assertion printed in the source. See
    % docs/paper-gaps/cpgsv17_periodic_sector_projector.tex.
    It remains to exclude nontrivial periodic vertical sectors under the
    literal canonical-form hypothesis. A cyclic projector $Q$ and an integer
    $p>1$ would commute with $[H^{(N)}]^p$ but not with $H^{(N)}$. Since
    $H^{(N)}$ is positive semidefinite, it is the positive $p$-th root of
    $[H^{(N)}]^p$; therefore every operator commuting with $[H^{(N)}]^p$
    also commutes with $H^{(N)}$, a contradiction. Together with
    invariant-projector closure, this exclusion must yield the normal vertical
    sectors and their spectral grouping.

    It also remains to construct the actual vertical representative groups
    and their corners over tensors $M_\alpha$. For each active copy
    $(\alpha,k)$, let $P_{\alpha,k}$ be its
    range projection. At some length its compression is nonzero and positive
    semidefinite, and its trace has the form
    \begin{align}
        \tr\!\left(P_{\alpha,k}H^{(N)}P_{\alpha,k}\right)
        &=\mu_{\alpha,k}d_\alpha^{(N)}>0.
        \label{eq:cpsv_prop_4_13_positive_weight}
    \end{align}
    Absorb the phase of one nonzero coefficient into $M_\alpha$, so that the
    reference coefficient satisfies $\mu_{\alpha,1}>0$. The identity for
    the original copy first gives $d_\alpha^{(N)}\ne0$. Hence the reference
    compression has nonzero trace $\mu_{\alpha,1}d_\alpha^{(N)}$ and is
    nonzero. Its positivity then gives $d_\alpha^{(N)}>0$, and
    (\ref{eq:cpsv_prop_4_13_positive_weight}) implies
    $\mu_{\alpha,k}>0$ for every copy.

    Once these actual corners are available, the pairwise Figure~8
    identity is complete under the literal canonical-form hypothesis.
    Theorem~\ref{thm:cpsv_literal_pairwise_vertical_gram_conjugation} gives
    \begin{align}
        (X_{\alpha,k}^\dagger X_{\alpha,k})M_\alpha^v
        (X_{\alpha,k}^\dagger X_{\alpha,k})^{-1}
        &=(X_{\alpha,1}^\dagger X_{\alpha,1})M_\alpha^v
        (X_{\alpha,1}^\dagger X_{\alpha,1})^{-1}.
        \label{eq:cpsv_prop_4_13_gram_conjugation}
    \end{align}
    If $M_\alpha$ is normal,
    Theorem~\ref{thm:cpsv_literal_positive_relative_gram_scalar} gives a real
    number $\omega_{\alpha,k}>0$ such that
    \begin{align}
        X_{\alpha,k}^\dagger X_{\alpha,k}
        &=\omega_{\alpha,k}
          X_{\alpha,1}^\dagger X_{\alpha,1}.
        \label{eq:cpsv_prop_4_13_gram}
    \end{align}
    These pairwise statements belong to Figures~7--8 in the proof of
    Proposition~4.13. They are distinct from the positive-fusion statement in
    Appendix~C.4 tracked by issue~\#4648.
    Choose $X_{\alpha,1}=\Id$ and replace each remaining gauge by
    $\omega_{\alpha,k}^{-1/2}X_{\alpha,k}$. The gauges are then unitary, while
    their conjugation actions are unchanged.

    Compose these unitary bond gauges with the mutually orthogonal inclusions
    of the retained vertical sectors, and take their adjoints as the block-row
    map $U$. The orthogonality of the sector ranges gives
    (\ref{eq:cpsv_prop_4_13_coisometry}), and the sector intertwinings give
    (\ref{eq:cpsv_prop_4_13_vertical_form}). Invariant-projector closure makes
    the retained support reducing, while the complement is the all-zero
    sector. Summing the sector reconstructions therefore gives the exact
    identity (\ref{eq:cpsv_prop_4_13_reconstruction}).

    The preceding paragraphs recast the source proof on the active retained
    coordinates required by the literal canonical-form definition.
\end{proof}

% Formalization status: the theorem above remains \notready. The theorem below
% proves the same conclusion from the stronger normalized BNT-refined horizontal
% form. For literal canonical form, the exact class-copy-plus-inactive
% coordinate regrouping, active representative sector decomposition, and
% retained-coordinate and original-coordinate forms of Lemma L, the original
% inserted-tensor equality, invariant-projector closure, and the conditional
% pairwise Figure-8/Gram identities are complete. Actual vertical representative
% groups and corners, literal periodic exclusion and spectral grouping, normalized
% grouped sectors, and the final coisometry remain open.

\begin{theorem}[BNT-refined horizontal form implies vertical canonical form]
    \label{thm:vertical_cf_of_horizontal_cf}
    \lean{MPOTensor.verticalCF_of_horizontalCF}
    \lean{MPOTensor.IsHorizontalCF.exists_cpsvVerticalDecomposition}
    \leanok
    \uses{def:mpdo, def:horizontal_cf_mpdo, def:vertical_cf_mpdo}
    A tensor in normalized BNT-refined horizontal form that generates matrix
    product density operators is also in canonical form in the vertical direction:
    there exist a basis of normal tensors $\{M_\alpha\}_\alpha$ for the
    vertically viewed tensor $\widetilde M$, positive diagonal matrices
    $\mu_\alpha$, and a coisometry $U$ from the physical space onto the
    retained nonzero sectors. Writing
    $B=\bigoplus_\alpha\mu_\alpha\otimes M_\alpha$, one has
    \begin{align}
        UU^\dagger&=\Id,
        \notag\\
        U\widetilde MU^\dagger&=B,
        \notag\\
        \widetilde M&=U^\dagger BU.
        \notag
    \end{align}
    The conclusion includes the basis-of-normal-tensors condition in
    \cite[Proposition~4.13]{Cirac2016MPDO_arXiv}, together with the positive
    weights and the coisometric realization. Its
    hypothesis is the normalized BNT-refined form of
    Definition~\ref{def:horizontal_cf_mpdo}, rather than the literal CPSV
    canonical form assumed in
    Theorem~\ref{thm:vertical_cf_of_cpsv_canonical_form}.
\end{theorem}

\begin{proof}
    \leanok
    \uses{def:mpv_phase_class_data,
        thm:mpdo_vertical_bnt_grouping_with_isometries,
        thm:mpdo_vertical_grouped_is_bnt,
        thm:mpdo_normalized_grouped_sector_maps,
        thm:vertical_cf_of_grouped_orthogonal_sectors}
    Apply
    Theorem~\ref{thm:mpdo_vertical_bnt_grouping_with_isometries} once and keep
    its representatives $M_\alpha$, multiplicities $r_\alpha$, exact
    positive weights $\mu_{\alpha,q}$, and physical sector maps. For these
    same witnesses,
    Theorem~\ref{thm:mpdo_vertical_grouped_is_bnt} proves that
    $\{M_\alpha\}_\alpha$ is a BNT for $\widetilde M$.

    Theorem~\ref{thm:mpdo_normalized_grouped_sector_maps} absorbs each
    normalized bond gauge into its physical map and gives maps
    $W_{\alpha,q}$ such that
    \begin{align}
        W_{\alpha,q}^\dagger W_{\alpha,q}&=\Id,
        \notag\\
        W_{\alpha,q}^\dagger W_{\beta,p}&=0
        \quad((\alpha,q)\ne(\beta,p)),
        \notag\\
        \widetilde M_vW_{\alpha,q}
        &=W_{\alpha,q}(\mu_{\alpha,q}M_\alpha^v),
        \notag\\
        \widetilde M_v
        &=\sum_\alpha\sum_q
          W_{\alpha,q}(\mu_{\alpha,q}M_\alpha^v)
          W_{\alpha,q}^\dagger.
        \notag
    \end{align}
    The canonical equivalence between the dependent disjoint union
    $\coprod_\alpha\{0,\ldots,r_\alpha-1\}$ and its standard finite index
    set reindexes the nested reconstruction sum without changing any weight
    or sector. The hypotheses of
    Theorem~\ref{thm:vertical_cf_of_grouped_orthogonal_sectors} now hold, and
    that theorem gives the vertical canonical form of $M$.
\end{proof}
